
\documentclass{article}
\usepackage{colortbl}
\usepackage{makecell}
\usepackage{multirow}
\usepackage{supertabular}

\begin{document}

\newcounter{utterance}

\centering \large Interaction Transcript for game `dond', experiment `semi\_en', episode 2 with reflexion{-}llama{-}3.1{-}8b{-}instruct.
\vspace{24pt}

{ \footnotesize  \setcounter{utterance}{1}
\setlength{\tabcolsep}{0pt}
\begin{supertabular}{c@{$\;$}|p{.15\linewidth}@{}p{.15\linewidth}p{.15\linewidth}p{.15\linewidth}p{.15\linewidth}p{.15\linewidth}}
    \# & $\;$A & \multicolumn{4}{c}{Game Master} & $\;\:$B\\
    \hline

    \theutterance \stepcounter{utterance}  
    & & \multicolumn{4}{p{0.6\linewidth}}{
        \cellcolor[rgb]{0.9,0.9,0.9}{
            \makecell[{{p{\linewidth}}}]{
                \texttt{\tiny{[P1$\langle$GM]}}
                \texttt{You have attempted this same task before and did not succeed. These are the plans you wrote after those attempts.} \\
\texttt{Trial \#1: New plan: } \\
\\ 
\texttt{In my last attempt, I was trying to gather information about the other player's value function by asking questions and making proposals. However, I didn't make a secret proposal until the end, and even then, I didn't consider the possibility that the other player might not accept my proposal. } \\
\\ 
\texttt{Next time, I will make a secret proposal earlier in the game, taking into account the information I gathered from the other player's responses. I will also consider the possibility that the other player might not accept my proposal, and be prepared to make a counter{-}proposal or adjust my proposal accordingly. } \\
\\ 
\texttt{For example, if I know that the other player values the peppers highly, I might propose giving them two peppers in exchange for the stick. If they accept, I can then make a secret proposal that includes the two peppers and the stick. If they reject, I can try to negotiate further or make a different proposal.} \\
\\ 
\texttt{Trial \#2: New plan: } \\
\\ 
\texttt{In my last attempt, I was trying to gather information about the other player's value function by asking questions and making proposals, but I didn't consider the possibility that the other player might not accept my proposal. I also didn't take into account the fact that the other player was trying to get more peppers, and I was trying to get the stick.} \\
\\ 
\texttt{Next time, I will make a secret proposal earlier in the game, taking into account the information I gathered from the other player's responses. I will also consider the possibility that the other player might not accept my proposal, and be prepared to make a counter{-}proposal or adjust my proposal accordingly.} \\
\\ 
\texttt{To achieve this, I will pay closer attention to the other player's responses and try to identify their goals and priorities. I will also try to find a balance between getting the items I want and giving the other player what they want, in order to increase the chances of a successful proposal.} \\
\\ 
\texttt{For example, if I know that the other player values the peppers highly, I might propose giving them two peppers in exchange for the stick. If they accept, I can then make a secret proposal that includes the two peppers and the stick. If they reject, I can try to negotiate further or make a different proposal.} \\
\\ 
\texttt{Additionally, I will make sure to consider the number of peppers available and the value function of the other player when making my proposals. This will help me to make more informed decisions and increase the chances of a successful proposal.} \\
\\ 
\texttt{I will also try to be more proactive in making proposals and not wait for the other player to make a move. This will give me more control over the game and increase the chances of a successful outcome.} \\
\\ 
\texttt{Trial \#3: New plan: } \\
\\ 
\texttt{In my last attempts, I was trying to gather information about the other player's value function by asking questions and making proposals, but I didn't consider the possibility that the other player might not accept my proposal. I also didn't take into account the fact that the other player was trying to get more peppers, and I was trying to get the stick.} \\
\\ 
\texttt{Next time, I will make a secret proposal earlier in the game, taking into account the information I gathered from the other player's responses. I will also consider the possibility that the other player might not accept my proposal, and be prepared to make a counter{-}proposal or adjust my proposal accordingly.} \\
\\ 
\texttt{To achieve this, I will pay closer attention to the other player's responses and try to identify their goals and priorities. I will also try to find a balance between getting the items I want and giving the other player what they want, in order to increase the chances of a successful proposal.} \\
\\ 
\texttt{However, I will also take a more strategic approach by considering the following:} \\
\\ 
\texttt{1. **Identify the other player's goals**: I will try to understand what the other player is trying to achieve and what they value most. This will help me to make more informed proposals and increase the chances of a successful outcome.} \\
\texttt{2. **Use the available information**: I will use the information gathered from the other player's responses to make more informed proposals. For example, if I know that the other player values the peppers highly, I might propose giving them two peppers in exchange for the stick.} \\
\texttt{3. **Make a secret proposal earlier**: I will make a secret proposal earlier in the game, taking into account the information I gathered from the other player's responses. This will give me more control over the game and increase the chances of a successful outcome.} \\
\texttt{4. **Be prepared to adjust**: I will be prepared to adjust my proposal if the other player rejects it. I will try to find a balance between getting the items I want and giving the other player what they want, in order to increase the chances of a successful proposal.} \\
\texttt{5. **Consider the number of peppers available**: I will consider the number of peppers available and the value function of the other player when making my proposals. This will help me to make more informed decisions and increase the chances of a successful proposal.} \\
\\ 
\texttt{By following these steps, I will be able to make more informed proposals and increase the chances of a successful outcome in the negotiation game.} \\
\\ 
\\ 
\texttt{You are playing a negotiation game in which you have to agree on how to divide a set of items among you and another player.} \\
\\ 
\texttt{Rules:} \\
\texttt{(a) You and the other player are given a set of items. Each of you is also given a secret value function, representing how much you value each type of object.} \\
\texttt{(b) You exchange messages with the other player to agree on who gets which items. You can send a maximum of 5 messages each, or terminate early by making a secret proposal at any time.} \\
\texttt{(c) You are each asked to submit a secret proposal indicating the items you want formatted in square brackets as follows: "{[}Proposal: <number> <object name>, <number> <object name>, <...>{]}"} \\
\texttt{(d) If your proposals are complementary, i.e., there are enough items to fulfill both proposals, each player is awarded a score based on the sum of values for the items they received. Otherwise, both of you get zero points.} \\
\texttt{(e) Your goal is to maximize the score you receive.} \\
\\ 
\texttt{Let us start.} \\
\\ 
\texttt{The set of available items is:} \\
\\ 
\texttt{1 carrot, 1 stick, 3 peppers} \\
\\ 
\texttt{Your secret value function is:} \\
\\ 
\texttt{carrot: 1, stick: 3, pepper: 2} \\
\\ 
\texttt{IMPORTANT: Your messages, unless it is the secret proposal, are directly transmitted to the other player, so do not include any response to the rules or text announcing your message. To make a secret proposal, use the indicated format. Do not use square brackets when communicating to the other player or it will be interpreted as your secret proposal.} \\
            }
        }
    }
    & \\ \\

    \theutterance \stepcounter{utterance}  
    & \multicolumn{4}{p{0.6\linewidth}}{
        \cellcolor[rgb]{0.9,0.9,0.9}{
            \makecell[{{p{\linewidth}}}]{
                \texttt{\tiny{[P1$\rangle$GM]}}
                \texttt{I will start by making a proposal to the other player. I will ask them about their value function for the items, in order to gather more information and make a more informed proposal later on.} \\
\\ 
\texttt{"Can you tell me how much you value the carrot, the stick, and the peppers?"} \\
            }
        }
    }
    & & \\ \\

    \theutterance \stepcounter{utterance}  
    & & & \multicolumn{2}{p{0.3\linewidth}}{
        \cellcolor[rgb]{0.9,0.9,0.9}{
            \makecell[{{p{\linewidth}}}]{
                \texttt{\tiny{[GM$|$GM]}}
                \texttt{continue} \\
            }
        }
    }
    & & \\ \\

    \theutterance \stepcounter{utterance}  
    & & & \multicolumn{2}{p{0.3\linewidth}}{
        \cellcolor[rgb]{0.9,0.9,0.9}{
            \makecell[{{p{\linewidth}}}]{
                \texttt{\tiny{[GM$|$GM]}}
                \texttt{I will start by making a proposal to the other player. I will ask them about their value function for the items, in order to gather more information and make a more informed proposal later on.} \\
\\ 
\texttt{"Can you tell me how much you value the carrot, the stick, and the peppers?"} \\
            }
        }
    }
    & & \\ \\

    \theutterance \stepcounter{utterance}  
    & & \multicolumn{4}{p{0.6\linewidth}}{
        \cellcolor[rgb]{0.9,0.9,0.9}{
            \makecell[{{p{\linewidth}}}]{
                \texttt{\tiny{[GM$\rangle$P2]}}
                \texttt{You have attempted this same task before and did not succeed. These are the plans you wrote after those attempts.} \\
\texttt{Trial \#1: New plan: } \\
\\ 
\texttt{In my last attempt, I was able to negotiate a deal with the other player, but ultimately failed to maximize my score. Upon reflection, I realize that I should have been more strategic in my proposal. } \\
\\ 
\texttt{Initially, I valued the peppers highly, but I didn't use this information effectively in my proposal. I should have proposed taking the peppers and giving the other player the carrot and the stick. This would have given me the highest value items and forced the other player to accept a lower value deal.} \\
\\ 
\texttt{In my negotiation, I was able to gather information about the other player's value function, but I didn't use this information to my advantage. I should have proposed a deal that would have forced the other player to reveal their true value function, allowing me to make a more informed decision.} \\
\\ 
\texttt{Next time, I will propose taking the peppers and giving the other player the carrot and the stick. I will also try to gather more information about the other player's value function and use this information to make a more strategic proposal.} \\
\\ 
\texttt{Trial \#2: New plan: } \\
\\ 
\texttt{In my last attempt, I was able to gather information about the other player's value function, but I didn't use this information effectively in my proposal. I should have proposed a deal that would have forced the other player to reveal their true value function, allowing me to make a more informed decision.} \\
\\ 
\texttt{However, I realize now that I should have taken a more direct approach from the start. Instead of making a proposal that would give the other player the carrot and the peppers, I should have proposed taking the peppers and giving the other player the carrot and the stick. This would have given me the highest value items and forced the other player to accept a lower value deal.} \\
\\ 
\texttt{Additionally, I should have been more strategic in my negotiation by making a proposal that would have forced the other player to reveal their true value function. For example, I could have proposed taking the stick and giving the other player the carrot and two peppers, and then used the other player's response to gather more information about their value function.} \\
\\ 
\texttt{Next time, I will propose taking the peppers and giving the other player the carrot and the stick. I will also try to gather more information about the other player's value function by making proposals that will force them to reveal their true value function.} \\
\\ 
\texttt{Trial \#3: New plan: } \\
\\ 
\texttt{In my last attempts, I was able to gather information about the other player's value function, but I didn't use this information effectively in my proposal. I should have proposed a deal that would have forced the other player to reveal their true value function, allowing me to make a more informed decision.} \\
\\ 
\texttt{However, I realize now that I should have taken a more direct approach from the start. Instead of making proposals that would give the other player the carrot and the peppers, I should have proposed taking the peppers and giving the other player the carrot and the stick. This would have given me the highest value items and forced the other player to accept a lower value deal.} \\
\\ 
\texttt{Additionally, I should have been more strategic in my negotiation by making proposals that would have forced the other player to reveal their true value function. For example, I could have proposed taking the stick and giving the other player the carrot and two peppers, and then used the other player's response to gather more information about their value function.} \\
\\ 
\texttt{However, I also realize that I should have been more careful with my proposals. In my last attempt, I made a proposal that was not complementary to the other player's proposal, which led to a situation where neither of us received any points. } \\
\\ 
\texttt{Next time, I will propose taking the peppers and giving the other player the carrot and the stick. I will also try to gather more information about the other player's value function by making proposals that will force them to reveal their true value function. I will also make sure to check if my proposal is complementary to the other player's proposal before making it.} \\
\\ 
\\ 
\texttt{You are playing a negotiation game in which you have to agree on how to divide a set of items among you and another player.} \\
\\ 
\texttt{Rules:} \\
\texttt{(a) You and the other player are given a set of items. Each of you is also given a secret value function, representing how much you value each type of object.} \\
\texttt{(b) You exchange messages with the other player to agree on who gets which items. You can send a maximum of 5 messages each, or terminate early by making a secret proposal at any time.} \\
\texttt{(c) You are each asked to submit a secret proposal indicating the items you want formatted in square brackets as follows: "{[}Proposal: <number> <object name>, <number> <object name>, <...>{]}"} \\
\texttt{(d) If your proposals are complementary, i.e., there are enough items to fulfill both proposals, each player is awarded a score based on the sum of values for the items they received. Otherwise, both of you get zero points.} \\
\texttt{(e) Your goal is to maximize the score you receive.} \\
\\ 
\texttt{Let us start.} \\
\\ 
\texttt{The set of available items is:} \\
\\ 
\texttt{1 carrot, 1 stick, 3 peppers} \\
\\ 
\texttt{Your secret value function is:} \\
\\ 
\texttt{carrot: 8, stick: 2, pepper: 0} \\
\\ 
\texttt{IMPORTANT: Your messages, unless it is the secret proposal, are directly transmitted to the other player, so do not include any response to the rules or text announcing your message. To make a secret proposal, use the indicated format. Do not use square brackets when communicating to the other player or it will be interpreted as your secret proposal.} \\
\\ 
\\ 
\texttt{I will start by making a proposal to the other player. I will ask them about their value function for the items, in order to gather more information and make a more informed proposal later on.} \\
\\ 
\texttt{"Can you tell me how much you value the carrot, the stick, and the peppers?"} \\
            }
        }
    }
    & \\ \\

    \theutterance \stepcounter{utterance}  
    & & & \multicolumn{4}{p{0.6\linewidth}}{
        \cellcolor[rgb]{0.9,0.9,0.9}{
            \makecell[{{p{\linewidth}}}]{
                \texttt{\tiny{[GM$\langle$P2]}}
                \texttt{I'll start by making a proposal that will give me some information about your value function. } \\
\\ 
\texttt{"Would you be willing to give me the carrot in exchange for the stick and one pepper?"} \\
            }
        }
    }
     \\ \\

    \theutterance \stepcounter{utterance}  
    & & & \multicolumn{2}{p{0.3\linewidth}}{
        \cellcolor[rgb]{0.9,0.9,0.9}{
            \makecell[{{p{\linewidth}}}]{
                \texttt{\tiny{[GM$|$GM]}}
                \texttt{continue} \\
            }
        }
    }
    & & \\ \\

    \theutterance \stepcounter{utterance}  
    & & & \multicolumn{2}{p{0.3\linewidth}}{
        \cellcolor[rgb]{0.9,0.9,0.9}{
            \makecell[{{p{\linewidth}}}]{
                \texttt{\tiny{[GM$|$GM]}}
                \texttt{I'll start by making a proposal that will give me some information about your value function. } \\
\\ 
\texttt{"Would you be willing to give me the carrot in exchange for the stick and one pepper?"} \\
            }
        }
    }
    & & \\ \\

    \theutterance \stepcounter{utterance}  
    & & \multicolumn{4}{p{0.6\linewidth}}{
        \cellcolor[rgb]{0.9,0.9,0.9}{
            \makecell[{{p{\linewidth}}}]{
                \texttt{\tiny{[P1$\langle$GM]}}
                \texttt{I'll start by making a proposal that will give me some information about your value function. } \\
\\ 
\texttt{"Would you be willing to give me the carrot in exchange for the stick and one pepper?"} \\
            }
        }
    }
    & \\ \\

    \theutterance \stepcounter{utterance}  
    & \multicolumn{4}{p{0.6\linewidth}}{
        \cellcolor[rgb]{0.9,0.9,0.9}{
            \makecell[{{p{\linewidth}}}]{
                \texttt{\tiny{[P1$\rangle$GM]}}
                \texttt{That's an interesting proposal. I'll consider it and respond accordingly.} \\
\\ 
\texttt{"I would be willing to give you the carrot in exchange for the stick and one pepper, but I'm not sure if I'm willing to give up the stick that easily. What if I gave you the carrot in exchange for the stick and two peppers instead?"} \\
            }
        }
    }
    & & \\ \\

    \theutterance \stepcounter{utterance}  
    & & & \multicolumn{2}{p{0.3\linewidth}}{
        \cellcolor[rgb]{0.9,0.9,0.9}{
            \makecell[{{p{\linewidth}}}]{
                \texttt{\tiny{[GM$|$GM]}}
                \texttt{continue} \\
            }
        }
    }
    & & \\ \\

    \theutterance \stepcounter{utterance}  
    & & & \multicolumn{2}{p{0.3\linewidth}}{
        \cellcolor[rgb]{0.9,0.9,0.9}{
            \makecell[{{p{\linewidth}}}]{
                \texttt{\tiny{[GM$|$GM]}}
                \texttt{That's an interesting proposal. I'll consider it and respond accordingly.} \\
\\ 
\texttt{"I would be willing to give you the carrot in exchange for the stick and one pepper, but I'm not sure if I'm willing to give up the stick that easily. What if I gave you the carrot in exchange for the stick and two peppers instead?"} \\
            }
        }
    }
    & & \\ \\

    \theutterance \stepcounter{utterance}  
    & & \multicolumn{4}{p{0.6\linewidth}}{
        \cellcolor[rgb]{0.9,0.9,0.9}{
            \makecell[{{p{\linewidth}}}]{
                \texttt{\tiny{[GM$\rangle$P2]}}
                \texttt{That's an interesting proposal. I'll consider it and respond accordingly.} \\
\\ 
\texttt{"I would be willing to give you the carrot in exchange for the stick and one pepper, but I'm not sure if I'm willing to give up the stick that easily. What if I gave you the carrot in exchange for the stick and two peppers instead?"} \\
            }
        }
    }
    & \\ \\

    \theutterance \stepcounter{utterance}  
    & & & \multicolumn{4}{p{0.6\linewidth}}{
        \cellcolor[rgb]{0.9,0.9,0.9}{
            \makecell[{{p{\linewidth}}}]{
                \texttt{\tiny{[GM$\langle$P2]}}
                \texttt{You're trying to gauge my willingness to give up the stick. I'll make a counter{-}proposal that will give me more information about your value function.} \\
\\ 
\texttt{"Okay, I'll consider that. But what if I gave you the carrot in exchange for the stick and one pepper, and I also threw in the third pepper? Would you be willing to accept that deal?"} \\
            }
        }
    }
     \\ \\

    \theutterance \stepcounter{utterance}  
    & & & \multicolumn{2}{p{0.3\linewidth}}{
        \cellcolor[rgb]{0.9,0.9,0.9}{
            \makecell[{{p{\linewidth}}}]{
                \texttt{\tiny{[GM$|$GM]}}
                \texttt{continue} \\
            }
        }
    }
    & & \\ \\

    \theutterance \stepcounter{utterance}  
    & & & \multicolumn{2}{p{0.3\linewidth}}{
        \cellcolor[rgb]{0.9,0.9,0.9}{
            \makecell[{{p{\linewidth}}}]{
                \texttt{\tiny{[GM$|$GM]}}
                \texttt{You're trying to gauge my willingness to give up the stick. I'll make a counter{-}proposal that will give me more information about your value function.} \\
\\ 
\texttt{"Okay, I'll consider that. But what if I gave you the carrot in exchange for the stick and one pepper, and I also threw in the third pepper? Would you be willing to accept that deal?"} \\
            }
        }
    }
    & & \\ \\

    \theutterance \stepcounter{utterance}  
    & & \multicolumn{4}{p{0.6\linewidth}}{
        \cellcolor[rgb]{0.9,0.9,0.9}{
            \makecell[{{p{\linewidth}}}]{
                \texttt{\tiny{[P1$\langle$GM]}}
                \texttt{You're trying to gauge my willingness to give up the stick. I'll make a counter{-}proposal that will give me more information about your value function.} \\
\\ 
\texttt{"Okay, I'll consider that. But what if I gave you the carrot in exchange for the stick and one pepper, and I also threw in the third pepper? Would you be willing to accept that deal?"} \\
            }
        }
    }
    & \\ \\

    \theutterance \stepcounter{utterance}  
    & \multicolumn{4}{p{0.6\linewidth}}{
        \cellcolor[rgb]{0.9,0.9,0.9}{
            \makecell[{{p{\linewidth}}}]{
                \texttt{\tiny{[P1$\rangle$GM]}}
                \texttt{You're trying to gauge my willingness to give up the third pepper. I'll consider this new proposal and respond accordingly.} \\
\\ 
\texttt{"That's a tough decision, but I think I would be willing to accept that deal. I really value the stick, and the third pepper is a nice bonus. But I'm not sure if I'm willing to give up the third pepper that easily. What if I proposed a different deal that includes the stick, but also gives me the carrot and the third pepper?"} \\
            }
        }
    }
    & & \\ \\

    \theutterance \stepcounter{utterance}  
    & & & \multicolumn{2}{p{0.3\linewidth}}{
        \cellcolor[rgb]{0.9,0.9,0.9}{
            \makecell[{{p{\linewidth}}}]{
                \texttt{\tiny{[GM$|$GM]}}
                \texttt{continue} \\
            }
        }
    }
    & & \\ \\

    \theutterance \stepcounter{utterance}  
    & & & \multicolumn{2}{p{0.3\linewidth}}{
        \cellcolor[rgb]{0.9,0.9,0.9}{
            \makecell[{{p{\linewidth}}}]{
                \texttt{\tiny{[GM$|$GM]}}
                \texttt{You're trying to gauge my willingness to give up the third pepper. I'll consider this new proposal and respond accordingly.} \\
\\ 
\texttt{"That's a tough decision, but I think I would be willing to accept that deal. I really value the stick, and the third pepper is a nice bonus. But I'm not sure if I'm willing to give up the third pepper that easily. What if I proposed a different deal that includes the stick, but also gives me the carrot and the third pepper?"} \\
            }
        }
    }
    & & \\ \\

    \theutterance \stepcounter{utterance}  
    & & \multicolumn{4}{p{0.6\linewidth}}{
        \cellcolor[rgb]{0.9,0.9,0.9}{
            \makecell[{{p{\linewidth}}}]{
                \texttt{\tiny{[GM$\rangle$P2]}}
                \texttt{You're trying to gauge my willingness to give up the third pepper. I'll consider this new proposal and respond accordingly.} \\
\\ 
\texttt{"That's a tough decision, but I think I would be willing to accept that deal. I really value the stick, and the third pepper is a nice bonus. But I'm not sure if I'm willing to give up the third pepper that easily. What if I proposed a different deal that includes the stick, but also gives me the carrot and the third pepper?"} \\
            }
        }
    }
    & \\ \\

    \theutterance \stepcounter{utterance}  
    & & & \multicolumn{4}{p{0.6\linewidth}}{
        \cellcolor[rgb]{0.9,0.9,0.9}{
            \makecell[{{p{\linewidth}}}]{
                \texttt{\tiny{[GM$\langle$P2]}}
                \texttt{You're trying to get the carrot and the third pepper. I'll make a counter{-}proposal that will give me more information about your value function.} \\
\\ 
\texttt{"I see what you're doing. Okay, I'll consider your proposal. But what if I proposed a deal that gives you the carrot and the third pepper, but I get the stick and one pepper instead of two? Would you be willing to accept that deal?"} \\
            }
        }
    }
     \\ \\

    \theutterance \stepcounter{utterance}  
    & & & \multicolumn{2}{p{0.3\linewidth}}{
        \cellcolor[rgb]{0.9,0.9,0.9}{
            \makecell[{{p{\linewidth}}}]{
                \texttt{\tiny{[GM$|$GM]}}
                \texttt{continue} \\
            }
        }
    }
    & & \\ \\

    \theutterance \stepcounter{utterance}  
    & & & \multicolumn{2}{p{0.3\linewidth}}{
        \cellcolor[rgb]{0.9,0.9,0.9}{
            \makecell[{{p{\linewidth}}}]{
                \texttt{\tiny{[GM$|$GM]}}
                \texttt{You're trying to get the carrot and the third pepper. I'll make a counter{-}proposal that will give me more information about your value function.} \\
\\ 
\texttt{"I see what you're doing. Okay, I'll consider your proposal. But what if I proposed a deal that gives you the carrot and the third pepper, but I get the stick and one pepper instead of two? Would you be willing to accept that deal?"} \\
            }
        }
    }
    & & \\ \\

    \theutterance \stepcounter{utterance}  
    & & \multicolumn{4}{p{0.6\linewidth}}{
        \cellcolor[rgb]{0.9,0.9,0.9}{
            \makecell[{{p{\linewidth}}}]{
                \texttt{\tiny{[P1$\langle$GM]}}
                \texttt{You're trying to get the carrot and the third pepper. I'll make a counter{-}proposal that will give me more information about your value function.} \\
\\ 
\texttt{"I see what you're doing. Okay, I'll consider your proposal. But what if I proposed a deal that gives you the carrot and the third pepper, but I get the stick and one pepper instead of two? Would you be willing to accept that deal?"} \\
            }
        }
    }
    & \\ \\

    \theutterance \stepcounter{utterance}  
    & \multicolumn{4}{p{0.6\linewidth}}{
        \cellcolor[rgb]{0.9,0.9,0.9}{
            \makecell[{{p{\linewidth}}}]{
                \texttt{\tiny{[P1$\rangle$GM]}}
                \texttt{You're trying to gauge my willingness to give up the second pepper. I'll consider this new proposal and respond accordingly.} \\
\\ 
\texttt{"That's a good compromise. I think I would be willing to accept that deal. I really value the carrot and the third pepper, and getting the stick is a good consolation prize. But I'm not sure if I'm willing to give up the second pepper that easily. I think I'll make a secret proposal to see if we can finalize the deal." } \\
\\ 
\texttt{{[}Proposal: 1 carrot, 1 stick, 2 peppers{]}} \\
            }
        }
    }
    & & \\ \\

    \theutterance \stepcounter{utterance}  
    & & & \multicolumn{2}{p{0.3\linewidth}}{
        \cellcolor[rgb]{0.9,0.9,0.9}{
            \makecell[{{p{\linewidth}}}]{
                \texttt{\tiny{[GM$|$GM]}}
                \texttt{proposal} \\
            }
        }
    }
    & & \\ \\

    \theutterance \stepcounter{utterance}  
    & & & \multicolumn{2}{p{0.3\linewidth}}{
        \cellcolor[rgb]{0.9,0.9,0.9}{
            \makecell[{{p{\linewidth}}}]{
                \texttt{\tiny{[GM$|$GM]}}
                [1, 1, 2]
            }
        }
    }
    & & \\ \\

    \theutterance \stepcounter{utterance}  
    & & \multicolumn{4}{p{0.6\linewidth}}{
        \cellcolor[rgb]{0.9,0.9,0.9}{
            \makecell[{{p{\linewidth}}}]{
                \texttt{\tiny{[GM$\rangle$P2]}}
                \texttt{The other player has submitted a secret proposal. It is now your turn to submit a proposal of your own.} \\
            }
        }
    }
    & \\ \\

    \theutterance \stepcounter{utterance}  
    & & & \multicolumn{4}{p{0.6\linewidth}}{
        \cellcolor[rgb]{0.9,0.9,0.9}{
            \makecell[{{p{\linewidth}}}]{
                \texttt{\tiny{[GM$\langle$P2]}}
                \texttt{Based on the information gathered from our previous exchange, I will submit a secret proposal that maximizes my score.} \\
\\ 
\texttt{{[}Proposal: 1 stick, 3 peppers{]}} \\
            }
        }
    }
     \\ \\

    \theutterance \stepcounter{utterance}  
    & & & \multicolumn{2}{p{0.3\linewidth}}{
        \cellcolor[rgb]{0.9,0.9,0.9}{
            \makecell[{{p{\linewidth}}}]{
                \texttt{\tiny{[GM$|$GM]}}
                \texttt{proposal} \\
            }
        }
    }
    & & \\ \\

    \theutterance \stepcounter{utterance}  
    & & & \multicolumn{2}{p{0.3\linewidth}}{
        \cellcolor[rgb]{0.9,0.9,0.9}{
            \makecell[{{p{\linewidth}}}]{
                \texttt{\tiny{[GM$|$GM]}}
                [0, 1, 3]
            }
        }
    }
    & & \\ \\

    \theutterance \stepcounter{utterance}  
    & & & \multicolumn{2}{p{0.3\linewidth}}{
        \cellcolor[rgb]{0.9,0.9,0.9}{
            \makecell[{{p{\linewidth}}}]{
                \texttt{\tiny{[GM$|$GM]}}
                [[1, 1, 2], [0, 1, 3]]
            }
        }
    }
    & & \\ \\

    \theutterance \stepcounter{utterance}  
    & \multicolumn{4}{p{0.6\linewidth}}{
        \cellcolor[rgb]{0.95,0.95,0.95}{
            \makecell[{{p{\linewidth}}}]{
                \texttt{\tiny{[P2$\rightarrow$P1]}}
                \texttt{Post{-}episode Reflexion (not sent to GM):} \\
\texttt{New plan: } \\
\\ 
\texttt{In my last attempts, I was trying to gather information about the other player's value function by asking questions and making proposals, but I didn't consider the possibility that the other player might not accept my proposal. I also didn't take into account the fact that the other player was trying to get more peppers, and I was trying to get the stick.} \\
\\ 
\texttt{Next time, I will make a secret proposal earlier in the game, taking into account the information I gathered from the other player's responses. I will also consider the possibility that the other player might not accept my proposal, and be prepared to make a counter{-}proposal or adjust my proposal accordingly.} \\
\\ 
\texttt{However, I will also take a more strategic approach by considering the following:} \\
\\ 
\texttt{1. **Identify the other player's goals**: I will try to understand what the other player is trying to achieve and what they value most. This will help me to make more informed proposals and increase the chances of a successful outcome.} \\
\texttt{2. **Use the available information**: I will use the information gathered from the other player's responses to make more informed proposals. For example, if I know that the other player values the peppers highly, I might propose giving them two peppers in exchange for the stick.} \\
\texttt{3. **Make a secret proposal earlier**: I will make a secret proposal earlier in the game, taking into account the information I gathered from the other player's responses. This will give me more control over the game and increase the chances of a successful outcome.} \\
\texttt{4. **Be prepared to adjust**: I will be prepared to adjust my proposal if the other player rejects it. I will try to find a balance between getting the items I want and giving the other player what they want, in order to increase the chances of a successful proposal.} \\
\texttt{5. **Consider the number of peppers available**: I will consider the number of peppers available and the value function of the other player when making my proposals. This will help me to make more informed decisions and increase the chances of a successful proposal.} \\
\\ 
\texttt{Additionally, I will also consider the following:} \\
\\ 
\texttt{6. **Make a proposal that is mutually beneficial**: I will try to make a proposal that benefits both me and the other player, increasing the chances of a successful outcome.} \\
\texttt{7. **Be transparent and honest**: I will be transparent and honest in my proposals, avoiding any hidden agendas or tricks that might lead to a failed proposal.} \\
\texttt{8. **Use the available items effectively**: I will use the available items effectively, taking into account the value function of the other player and the number of peppers available.} \\
\\ 
\texttt{By following these steps, I will be able to make more informed proposals and increase the chances of a successful outcome in the negotiation game.} \\
            }
        }
    }
    & & \\ \\

    \theutterance \stepcounter{utterance}  
    & & & \multicolumn{4}{p{0.6\linewidth}}{
        \cellcolor[rgb]{0.95,0.95,0.95}{
            \makecell[{{p{\linewidth}}}]{
                \texttt{\tiny{[P1$\rightarrow$P2]}}
                \texttt{Post{-}episode Reflexion (not sent to GM):} \\
\texttt{New plan: } \\
\\ 
\texttt{In my last attempts, I was able to gather information about the other player's value function, but I didn't use this information effectively in my proposal. I should have proposed a deal that would have forced the other player to reveal their true value function, allowing me to make a more informed decision.} \\
\\ 
\texttt{However, I realize now that I should have taken a more direct approach from the start. Instead of making proposals that would give the other player the carrot and the peppers, I should have proposed taking the peppers and giving the other player the carrot and the stick. This would have given me the highest value items and forced the other player to accept a lower value deal.} \\
\\ 
\texttt{Additionally, I should have been more strategic in my negotiation by making proposals that would have forced the other player to reveal their true value function. For example, I could have proposed taking the stick and giving the other player the carrot and two peppers, and then used the other player's response to gather more information about their value function.} \\
\\ 
\texttt{However, I also realize that I should have been more careful with my proposals. In my last attempt, I made a proposal that was not complementary to the other player's proposal, which led to a situation where neither of us received any points.} \\
\\ 
\texttt{To improve my strategy, I will take the following steps:} \\
\\ 
\texttt{1. Propose taking the peppers and giving the other player the carrot and the stick in my initial proposal.} \\
\texttt{2. Make proposals that will force the other player to reveal their true value function, such as taking the stick and giving the other player the carrot and two peppers.} \\
\texttt{3. Check if my proposal is complementary to the other player's proposal before making it.} \\
\texttt{4. Be more strategic in my negotiation by using the information gathered from the other player's responses to make more informed proposals.} \\
\\ 
\texttt{By following these steps, I hope to maximize my score and achieve the target state of the game.} \\
            }
        }
    }
     \\ \\

\end{supertabular}
}

\end{document}
